\subsection{Kiểm thử hệ thống}

Mục này trình bày môi trường thử nghiệm và kết quả của ba nhóm kiểm thử được thực hiện trực tiếp trên hệ thống đang chạy: kiểm thử tải (P1--P4), kiểm thử toàn vẹn (P5--P6) và thống kê mạng blockchain Hyperledger Fabric.

\subsubsection{Môi trường triển khai}

Toàn bộ hạ tầng -- bao gồm mạng Hyperledger Fabric (3 orderer, 2 peer), các vi dịch vụ NestJS, MongoDB (replica set), Redis -- được chạy trên cùng một máy vật lý thông qua Docker Compose với cơ sở dữ liệu cô lập (\textit{*\_test}) để tránh ảnh hưởng dữ liệu sản xuất. Cấu hình phần cứng và phần mềm được liệt kê tại Bảng~\ref{tab:hw-env} và Bảng~\ref{tab:sw-env}.

\begin{table}[H]
    \centering
    \caption{Cấu hình phần cứng môi trường thử nghiệm}
    \label{tab:hw-env}
    \begin{tabularx}{\textwidth}{|L{0.35}|L{0.65}|}
        \hline
        \textbf{Thành phần} & \textbf{Thông số} \\
        \hline
        Thiết bị  & Dell Vostro 5620 \\
        \hline
        CPU       & Intel Core i5-1240P (12 nhân, 16 luồng, tối đa 4{,}4~GHz) \\
        \hline
        RAM       & 24~GB DDR4 \\
        \hline
        Ổ cứng    & 512~GB NVMe SSD \\
        \hline
        GPU       & Intel Iris Xe Graphics \\
        \hline
        Kiến trúc & x86\_64 (64-bit) \\
        \hline
    \end{tabularx}
\end{table}

\begin{table}[H]
    \centering
    \caption{Cấu hình phần mềm môi trường thử nghiệm}
    \label{tab:sw-env}
    \begin{tabularx}{\textwidth}{|L{0.35}|L{0.65}|}
        \hline
        \textbf{Thành phần} & \textbf{Phiên bản / Giá trị} \\
        \hline
        Hệ điều hành            & Fedora Linux 43 Workstation Edition \\
        \hline
        Kernel                  & Linux 7.0.10 \\
        \hline
        Kiến trúc hệ điều hành  & x86\_64 \\
        \hline
        Ảo hóa                  & Intel VT-x \\
        \hline
    \end{tabularx}
\end{table}

\subsubsection{Kết quả kiểm thử tải}

Bốn pha kiểm thử tải lần lượt mô phỏng 50, 100, 500 và 1.000 cử tri thực hiện đầy đủ luồng bỏ phiếu -- đăng nhập, khởi tạo phiên, ký mù ngưỡng, nộp phiếu lên blockchain, đóng phiếu, tiết lộ và kiểm phiếu -- với độ trễ ngẫu nhiên trước mỗi yêu cầu nhằm tái hiện mô hình phân tán của một cuộc bầu cử thực tế. Bảng~\ref{tab:load-summary} tổng hợp kết quả.

\begin{table}[H]
    \centering
    \caption{Tổng hợp kết quả kiểm thử tải (P1--P4)}
    \label{tab:load-summary}
    \small
    \begin{tabularx}{\textwidth}{|C{1}|C{1}|C{1}|C{1}|C{1}|C{1}|C{1}|C{1}|}
        \hline
        \multirow{2}{*}{\textbf{Pha}} &
        \multirow{2}{*}{\textbf{Cử tri}} &
        \multirow{2}{*}{\textbf{OK}} &
        \multirow{2}{*}{\textbf{Lỗi}} &
        \multicolumn{2}{c|}{\textbf{VOTE (ms)}} &
        \multicolumn{2}{c|}{\textbf{REVEAL (ms)}} \\
        \cline{5-8}
        & & & & p50 & p95 & p50 & p95 \\
        \hline
        P1 &  50 &  50 & 0 & 2\,208 & 2\,260 & 2\,097 & 2\,238 \\
        \hline
        P2 & 100 & 100 & 0 & 2\,201 & 2\,263 & 2\,084 & 2\,117 \\
        \hline
        P3 & 500 & 500 & 0 &   755  & 1\,571 &   420  & 1\,296 \\
        \hline
        P4 & 1\,000 & 1\,000 & 0 &   638  & 1\,742 &   381  & 1\,384 \\
        \hline
    \end{tabularx}
\end{table}

Cả bốn pha đều hoàn thành không có lỗi. Bảng~\ref{tab:vote-substeps} phân tách độ trễ từng bước con trong pha VOTE của P1 và P3, làm rõ vị trí nút cổ chai và ảnh hưởng của cơ chế batch.

\begin{table}[H]
    \centering
    \caption{Độ trễ từng bước con của pha VOTE, so sánh P1 và P3 (đơn vị: ms)}
    \label{tab:vote-substeps}
    \small
    \begin{tabularx}{\textwidth}{|L{2}|C{0.75}|C{0.75}|C{0.75}|C{0.75}|}
        \hline
        \multirow{2}{*}{\textbf{Bước}} &
        \multicolumn{2}{c|}{\textbf{P1 (50 cử tri)}} &
        \multicolumn{2}{c|}{\textbf{P3 (500 cử tri)}} \\
        \cline{2-5}
        & p50 & p95 & p50 & p95 \\
        \hline
        sign-in              &  75    &  89    & 116   & 280   \\
        \hline
        start-session        &  22    &  26    &  43   & 100   \\
        \hline
        sign (threshold)     &  15    &  19    &  26   &  68   \\
        \hline
        \textbf{submit (blockchain)} & \textbf{2\,088} & \textbf{2\,132} & \textbf{487} & \textbf{1\,352} \\
        \hline
    \end{tabularx}
\end{table}

Ba bước đầu (sign-in, start-session, threshold sign) có độ trễ thấp và ổn định (15--116~ms). Toàn bộ \textbf{nút cổ chai} tập trung ở bước \textit{submit (blockchain)}: p50 = 2\,088~ms ở P1 so với 487~ms ở P3. Sự chênh lệch này phản ánh cơ chế batch của orderer -- ở tải thấp (P1/P2), block bị cắt sớm theo \textit{BatchTimeout} khoảng 2s trước khi đầy giao dịch; ở tải cao (P3/P4), các giao dịch đến đủ nhanh để lấp đầy block trước khi hết timeout, rút thời gian chờ xuống đáng kể. Ở P4 (1.000 cử tri), p50 của bước submit tiếp tục giảm xuống còn khoảng 410~ms nhờ mật độ giao dịch cao hơn, trong khi p95 tăng nhẹ do hàng đợi tại tầng endorsement.

\subsubsection{Kết quả kiểm thử toàn vẹn}

Kiểm thử toàn vẹn xác minh hệ thống từ chối tám lớp vi phạm nghiêm trọng trong điều kiện vận hành thực. Bảng~\ref{tab:negative-tests} trình bày kết quả trên hai pha P5 (100 cử tri, tải phân bố đều theo thời gian) và P6 (500 cử tri, tải tập trung trong thời gian ngắn).

\begin{table}[H]
    \centering
    \caption{Kết quả kiểm thử toàn vẹn trên pha P5 và P6}
    \label{tab:negative-tests}
    \small
    \begin{tabular*}{\textwidth}{@{\extracolsep{\fill}}|p{5.5cm}|p{4.5cm}|c|c|}
        \hline
        \textbf{Kịch bản vi phạm} & \textbf{Phản hồi kỳ vọng} & \textbf{P5} & \textbf{P6} \\
        \hline
        Bỏ phiếu hai lần (double-vote)                  & 409 Already voted             & PASS & PASS \\
        \hline
        Nộp phiếu chữ ký giả mạo                        & 409 Signature mismatch        & PASS & PASS \\
        \hline
        Bỏ phiếu sau khi bầu cử đóng                    & 400 Election not active       & PASS & PASS \\
        \hline
        Tiết lộ phiếu chữ ký giả mạo                    & 400 Invalid signature         & PASS & PASS \\
        \hline
        Tiết lộ ứng viên ngoài danh sách                & 403 Candidate not in election & PASS & PASS \\
        \hline
        Tiết lộ phiếu trùng lặp (double-reveal)         & 409 Already revealed          & PASS & PASS \\
        \hline
        Replay nonce ký mù đã sử dụng                   & 409 Nonce already used        & PASS & PASS \\
        \hline
        Truy xuất API quản trị không có quyền            & 401 Unauthorized              & PASS & PASS \\
        \hline
        \multicolumn{2}{|l|}{\textbf{Tỉ lệ vượt qua}} & \textbf{8/8} & \textbf{8/8} \\
        \hline
    \end{tabular*}
\end{table}

Toàn bộ tám kịch bản vi phạm đều bị từ chối đúng như thiết kế ở cả hai pha. Cơ chế kiểm tra trạng thái phiên Redis kết hợp với chỉ mục duy nhất MongoDB trên cặp \pkg{(electionId, voterId)} đảm bảo mọi nỗ lực bỏ phiếu hoặc tiết lộ trùng lặp (\textit{double-vote}, \textit{double-reveal}) đều bị phát hiện và từ chối ngay tại tầng dịch vụ, trước khi yêu cầu đến tầng blockchain. Tấn công replay nonce bị chặn nhờ danh sách đen nonce phía server; truy xuất trái phép API quản trị bị từ chối bởi middleware xác thực JWT hai cặp khóa.

\subsubsection{Thống kê mạng blockchain Hyperledger Fabric}

Bảng~\ref{tab:block-stats} tổng hợp chỉ số sổ cái qua sáu pha kiểm thử (tổng 1.137 block, 4.506 giao dịch).

\begin{table}[H]
    \centering
    \caption{Thống kê mạng blockchain Hyperledger Fabric theo pha kiểm thử}
    \label{tab:block-stats}
    \small
    \begin{tabularx}{\textwidth}{|C{1}|C{1}|C{1}|C{1}|C{1}|C{1}|}
        \hline
        \textbf{Pha} & \textbf{Cử tri} & \textbf{\#tx} &
        \textbf{tx/block (tb)} & \textbf{Interval median (s)} & \textbf{Throughput (tx/s)} \\
        \hline
        P1 &    50 &  101 & 1{,}28 & 6{,}33 & 0{,}17 \\
        \hline
        P2 &   100 &  201 & 1{,}69 & 3{,}93 & 0{,}34 \\
        \hline
        P3 &   500 & 1001 & 4{,}96 & 0{,}87 & 4{,}98 \\
        \hline
        P4 & 1\,000 & 2001 & 4{,}97 & 0{,}72 & 6{,}90 \\
        \hline
        P5 &   100 &  201 & 1{,}65 & 4{,}10 & 0{,}34 \\
        \hline
        P6 &   500 & 1001 & 4{,}72 & 0{,}93 & 4{,}61 \\
        \hline
    \end{tabularx}
\end{table}

Một số nhận xét rút ra từ dữ liệu sổ cái:
\begin{itemize}
    \item \textbf{Mô hình commit--reveal nhất quán:} Tỉ lệ giao dịch vote:reveal $\approx$ 1:1 ổn định qua mọi pha; mỗi cuộc bầu cử có đúng một giao dịch gốc Merkle duy nhất tại thời điểm đóng phiếu, xác nhận tính đầy đủ của dữ liệu trên sổ cái.
    \item \textbf{Hành vi batch của orderer:} Ở tải thấp (P1/P2/P5), phần lớn block chỉ chứa 1--2 giao dịch do block bị cắt theo BatchTimeout khoảng 2s trước khi đủ số lượng tối đa; ở tải cao (P3/P4), các giao dịch đến đủ nhanh để lấp đầy block trước khi hết timeout -- interval median giảm từ 6{,}33s (P1) xuống 0{,}87s (P3) và 0{,}72s (P4), throughput tăng tương ứng lên 6{,}90~tx/s.
    \item \textbf{Khả năng mở rộng từ 500 lên 1.000 cử tri:} Khi tải nhân đôi từ P3 lên P4, throughput tăng từ 4{,}98~tx/s lên 6{,}90~tx/s; p50 của bước VOTE giảm từ 755~ms xuống 638~ms nhờ mật độ giao dịch cao hơn lấp đầy block sớm hơn, trong khi p95 tăng nhẹ từ 1{,}571~ms lên 1{,}742~ms do hàng đợi tại tầng peer endorsement -- hệ thống vẫn vận hành ổn định, không có lỗi.
\end{itemize}
